\documentclass{mjsart}
\title{天津科技大学开放原子开源协会章程}
\begin{document}
\maketitle
\section{总则}

\textbf{第一条：协会性质}

天津科技大学开放原子开源协会隶属于天津科技大学人工智能学院，是为校内热爱开源技术、愿意投身开源实践与学术创新的学生创建的学术科技类学生社团组织，接受院团委直接领导和统一管理，需在遵守学校有关规定的前提下开展活动。

\textbf{第二条：协会目的}

本协会致力于推动开源技术在校园内的普及与实践，搭建开源技术交流与实践平台，宣传开源文化，培养学生的开源社区贡献能力，助力学生提升技术应用与学术创新素养，促进学术创新与技术应用融合发展，为相关领域人才培养提供支持。

\textbf{第三条：协会宗旨}

坚守开源共享、协作创新的理念，帮助同学们深入学习开源技术，激发技术探索与创新热情，实现理论知识与开源实践的深度结合，营造积极向上的校园开源文化氛围。

\section{组织与管理}

\textbf{第四条：社团活动发起条件}

（一）有 2 名以上会员倡导发起；

（二）有专门负责人统筹筹备活动相关事宜；

（三）完成完整准备工作，包括明确活动名称、宗 旨、主要任务、活动内容、经费来源、物质条件及其他应说明的事项。

\textbf{第五条：活动准则}

社团活动必须奉行公开透明原则，发布广告、公告等材料须注明协会名称，不得盗用指导部门或其他组织的名义开展活动。协会举办各类活动前，须提前一周向协会相关负责部门提交申请报告、经费预算及完整活动方案，经学院相关部门批准后方可实施。活动结束后，需递交活动总结一份，在协会内部存档备案。

\textbf{第六条：联合活动管理}

协会与校内其他单位或校外团体、单位举办联合活动，须事先征得学院同意，并上报合作单位的证明材料及活动方案，经学院研究批准后方可进行。联合活动结束后，需将书面总结材料及音像资料交至相关部门存档；学期末应将本学期的活动总结及财务收支状况上报学院相关管理部门。

\textbf{第七条：涉外活动管理}

协会举办有外籍专家、校外机构人员参与的活动，必须上报学院批准后方可进行。

\textbf{第八条：活动范围限制}

协会不得开展任何以盈利为目的的商业性活动，也不得开展超出其宗旨的活动。

\textbf{第九条：讲座、报告类活动规定}

（一） 协会邀请社外人员在校举办讲座、报告等室内活动，须提前一周向协会秘书处相关岗位提出申请。 申请时需说明被邀请人的姓名、所在单位、讲座或报告内容、时间、地点、主要负责人等信息，必要时应出示被邀请人的身份证明或单位介绍信。

（二） 经协会内部审核并上报学院审核同意后，活动负责人持相关证明到学校有关部门借用活动场所。

（三） 讲座、报告会的秩序由协会负责维持，若讲座、报告会的内容与申请不符，一切后果由协会及组织者承担。

\textbf{第十条：活动责任划分}

协会活动责任由协会自行承担。协会负责人及其主要成员在活动过程中有错误或重大失误的，个人也应承担相应责任。协会负责人或其主要成员以协会名义所进行的活动视为协会活动。

\textbf{第十一条：奖惩机制}

每学年协会根据活动开展情况，评选优秀协会活动和先进个人并给予奖励。奖励方式包括通报表扬、颁发奖状、证书、奖章、奖品或奖金等。若协会严重违反相关规定，学院将视情节对协会负责人及其他负有责任的成员给予批评教育；对有严重错误者， 由相关部门上报学生工作办公室予以校纪处分。因在协会活动中违反规定受到校纪处分的学生，未经学院同意，不得再参加任何协会活动。协会会员有权了解所在协会的章程、组织机构和财务制度，可对协会的管理和活动提出建议，也可向学院相关部门投诉或申诉。

\textbf{第十二条：工作会议制度}

为加强对协会工作的指导和管理，定期召开协会工作会议，协会成员需按时参加。

\section{成员}

\textbf{第十三条：入社条件}

天津科技大学学生、研究生认可本章程，在百团纳新或其他协会管理中心批准的时间经考核通过，即可加入本协会。活跃会员大会可规定其他不与学校规定相抵触的加入方式。 同时，要求成员对开源技术有一定兴趣，具备基本的学习能力和协作意识。

\textbf{第十四条：会员权利}

（一） 参与协会组织的各类活动；

（二） 对协会工作提出建议与批评；

（三） 符合活跃会员标准时行使表决权、选举权及被选举权；

（四） 有权了解协会的章程、组织机构和财务制度，对协会的管理和活动提出质询；

（五） 有权按照本章程自由加入或退出本协会。

\textbf{第十五条：会员义务}

（一） 遵守章程及协会管理制度；

（二） 积极参与协会活动，履行相应职责；

（三） 维护协会声誉与合法权益；

（四） 接受协会的定期注册。

\textbf{第十六条：活跃会员认定}

（一） 过去六个月内参与协会正式活动次数占比大于等于

25\%者、过去六个月内参与为期长于一个月的长期活动一次以上者和过去六个月继续参与六个月前开启的活动者，为活跃会员。

（二） 新加入未满六个月但已满 30 日的会员，加入后参与协会正式活动次数占比大于等于25\%的、加入后参与为期长于一个月的长期活动一次以上的和加入后参与六个月前开启的长期活动的，为活跃会员。

（三） 新加入未满 30 日的会员，不认定为活跃会员。

（四） 本协会的创始会员， 自动成为活跃会员，且不受前款所述对于新加入会员的限制；但在本协会第一次注册前参加活动少于二次者除外。

（五） 活跃会员在过去六个月参与活动比例低于25\%且未参与为期长于一个月的长期活动一次以上的，其活跃会员资格自然丧失。

（六） 活跃会员资格每月认定一次，在该月第一次活跃会员大会召开前宣布，活跃会员名单由主席团公告。只有活跃会员可以出席活跃会员大会并行使表决权，非活跃会员可以自愿列席。

\section{组织机构与职责}

\textbf{第十七条：活跃会员大会}

（一）性质：活跃会员大会是协会的决策机关。

（二）职权：

1. 决定协会发展规划及重大事项；

2. 审查和批准协会预算（含开放原子基金会支持积分）；

3. 经三分之二以上活跃会员出席，并经出席的活跃会员的三分之二以上多数赞成，可修改本章程、开除会员；

4. 选举主席团；

5. 根据主席团的提名，决定财务专员的人选；

6. 当主席团成员或财务专员出现严重行为不端时，罢免主席团成员或财务专员。罢免主席团成员职务，必须经过学校规定的程序进行；罢免财务专员，需要经过指导老师同意。

（三）议事规则：

1. 决议可通过线上投票方式作出；

2. 举行选举时，有投票权的会员有权了解候选人情况、要求改变候选人、不选任何一个候选人或者另选他人；

3. 举行其他表决时，有投票权的会员有权了解草案内容和相关情况，有权要求修改草案，有权投赞成、反对和弃权票；

4. 除本章程另有规定外，举行选举和表决须有半数以上有投票权的会员参与，并以出席人员的半数以上赞成票作出决定。

\textbf{第十八条：主席团}

（一） 产生与任期： 由活跃会员大会选举产生候选人，并按照学院有关规定任命，任期一年。主席团主席由主席团内部选举产生，召集并主持会议；主席的主席团成员职务被罢免时，其主席职务自动终止。主席团成员任满一届以上后，继续担任主席团成员应当根据学校有关规定通过审批；主席团成员未任满一届的，视为任满一届任期。

（二） 职责：

1. 召集并主持活跃会员大会，预先讨论拟提请活跃会员大会讨论的问题，决定是否列入议程（网络投票除外）；

2. 执行活跃会员大会决议；

3. 策划和组织日常协会活动；

4. 提名财务专员人选；

5. 行使活跃会员大会决定的其他职权，以及为完成学校有关部门交办的任务所需要的职权。

\textbf{第十九条：财务专员}

（一） 设立依据：按照学校有关规定，本协会设立财务专员。

（二） 职权：

1. 保管协会全部财产；

2. 按协会预算和主席团决议支配资金和开放原子基金会资助；

3. 每季度公示资产使用明细（含基金会资助）。

（三）参会规定：财务专员不是主席团成员时，列席主席团会议。

\textbf{第二十条：部门设置及职责}

（一）秘书处：负责协会日常行政统筹与协调工作，承担会议组织、议程拟定及会议记录归档；负责活动场地、设备的申请与统筹安排，做好活动签到、物料管理等后勤保障；制定协会内部管理制度，统筹各类文件、资料的整理、排版与存档，协调各部门间工作衔接；承担外联相关工作，负责对接校内校外其他协会、学院相关部门及校外开源机构、合作单位，拓展合作资源与交流渠道，保障协会对外协作活动顺利开展，全面保障协会运营高效有序。

（二） 项目部：负责协会各类开源相关活动的策划与落地执行，包括技术沙龙、项目实践营、开源竞赛等；牵头对接开源社区、合作企业的项目资源，组织会员参与开源项目贡献；负责活动方案的细化设计、流程把控、现场组织及活动总结复盘，收集活动反馈并优化后续活动策划。

（三） 宣传部：承担协会品牌形象塑造与宣传推广工作，负责活动海报、宣传册等物料设计制作；运营协会官方新媒体平台（微信公众号、视频号等），撰写活动新闻稿、科普推文，发布活动预告与回顾；负责活动现场拍摄、影像资料整理归档；组织内部宣传技能培训，提升宣传内容质量与传播效果，扩大协会校园影响力。

（四） 运维部：负责协会技术基础设施的搭建与维护，包括协会官网、线上交流平台、项目代码仓库等的日常运维；为协会活动提供技术支持，保障线上活动直播、线下活动设备调试等顺利开展；协助技术实践类活动的环境搭建, 指导会员掌握基础运维技能与开源工具使用方法。

（五） 学术部：负责协会学术交流与知识普及工作，组织开源技术讲座、专题研讨、文献分享会等学术活动；整理开源技术学习资料，编制学习手册，搭建会员学习互助平台；对接校内专业教师、校外开源技术专家，为会员提供学术指导与技术答疑；跟踪开源领域前沿动态，梳理行业发展趋势并分享给会员，助力提升会员学术素养与技术认知。

\section{财务管理}

\textbf{第二十一条：经费来源}

（一） 学校及学院专项资助；

（二） 开放原子基金会等社会支持；

（三） 其他合法和合乎学校规定的收入。

\textbf{第二十二条：经费管理原则}

协会经费遵循“专款专用、合理节约 ”的原则，用于协会的管理及各项活动组织，不得用于与协会活动无关的方面。协会所有财产属协会全体成员共有，必须用于服务会员、发展协会。

\textbf{第二十三条：经费使用与报销规定}

（一） 购买日常办公用品，需逐项列出购买用品清单，经主席团批准同意后方可购买。金额低于主席团规定的一定数额时，可由秘书长签字批准；金额高于主席团规定的数额时，必须由主席团批准。

（二） 购买后， 由秘书处点数目核对无误后，方可凭收据和发票给予报销。

（三） 凡发票缺损、涂改、丢失或发票不符合填写标准者，不予报销；凡未经批准自行决定购置者，不予报销。

（四）举办大型活动时，组织策划者须预算整个活动的经费（如宣传、联系、评委、奖品、场地等费用），逐项列清后交至秘书处审批；

（五）秘书处应当在每笔开销发生的第二 日起三日内，将该笔开支及相关材料移送财务专员保存。

\textbf{第二十四条：财务公开}

每学期第十七周为协会财务公开日， 由财务专员向全体会员公示本学期财务收支状况，接受会员监督。

\section{换届与纳新}

\textbf{第二十五条：换届原则}

参与竞聘职务涉及协会主席及各部长、副部长，采取按需设岗、能者居之的原则，可连选连任，重视和尊重人才。所有竞聘同学，须服从换届委员会协商后的岗位协调，优化组织内部结构。

\textbf{第二十六条：部长竞聘条件}

（一） 热爱社会主义祖国，拥护党的基本路线，有正确的政治立场、观点和态度，具有良好的品行、修养和道德风尚；

（二） 遵纪守法，模范执行《高校学生行为准则》和学校的各项规章制度， 自觉参加校园文明建设活动，能够在各方面起模范带头作用；

（三） 学习目的明确，态度端正，刻苦勤奋，学有余力，成绩良好，在同学中有较高的威信；

（四） 积极参加社会实践、文化、科技活动和公益活动，集体观念强，具有较强的组织协调能力和创新意识，组织观念强；

（五） 甘愿奉献，热心学生工作，并希望在工作中锻炼提高自己的能力；

（六） 在校期间无违反校纪校规或受到校纪校规处分的情况；

（七） 学生党员干部、学习成绩优秀者、在其他学生组织机构无兼职者优先；

（八） 在校期间无违反校纪校规或受到校纪校规处分的情况；

（九） 未因违反有关规定被撤职或在曾负责的协会被宣布解散时承担主要责任。

\textbf{第二十七条：部长换届流程}

（一） 报名：竞聘采取自愿报名的方式。符合竞聘资格且有意参加竞聘的同学，于规定时间内上交相关材料；

（二） 资格审查： 由秘书处审查资格并确定竞职演说人员；

（三） 竞聘演讲：竞职人员在规定时间和地点进行竞聘演说， 内容包括竞聘意向、在校主要表现、对所应聘部门的理解及具体工作设想等。台下观众可提问，竞选者进行作答。主席团人员根据演讲情况投票，确定入围人选。竞聘演说采取同岗位回避制度；

（四） 投票确定候选人：主席团成员进行全体投票，确定各岗位候选人；

（五） 组织考察：主席团对入围人员进行全面考察，考察内容包括德、绩、能、勤等。考察方式包括到班级了解情况、找本人谈话，考察组写出考察意见；

（六） 投票任命：活跃会员大会投票决定是否任命。活跃会员大会决定部长人选应有半数以上的活跃会员出席，以出席活跃会员的半数以上赞成通过；

（六）张榜公示：于任命后3 日内统一张榜公示。新一届部长试用期三个月（从张榜公示之日起算）。

\textbf{第二十八条：纳新流程}

任何天津科技大学学生、研究生在百团纳新期间提出申请并参加培训，面试通过者即可加入本协会。

\textbf{第二十九条：主席换届流程}

（一） 报名：采取自愿报名的方式。符合报名资格且有意参加选举的同学，于规定时间内如实填写并上交相关材料。报名竞选主席的人员可申请增报一个部长职位；

（二） 资格审查： 由秘书处审查资格并确定竞职演说人员；

（三） 竞聘演讲：竞职人员在活跃会员大会上进行竞聘演说， 内容包括竞聘意向、在校主要表现、对所应聘部门的理解及具体工作设想等。全体会员可提问，竞选者进行作答。活跃会员大会根据演讲情况投票，确定入围人选。竞聘演说采取同岗位回避制度；

（四） 报请学院根据学院有关规定确定最终人选。

\section{协会负责人（主席团成员）}

\textbf{第三十条：额外任职条件}

（一） 在思想上，有良好的道德情操，正确处理个人利益和集体利益的关系，克己奉公，助人为乐，团结同学，诚实谦虚，勇于批评和自我批评，严于律己，宽以待人；

（二） 在政治上，具有坚定正确的方向和马克思主义基础知识，能正确理解党的路线、方针和政策，能用马克思主义的观点分析问题， 自觉维护和投身改革，积极参加社会政治活动，具有强烈的事业心和责任心，有全心全意为会员服务的献身精神；

（三） 在工作上，有一定的思维判别能力和决策能力，能熟练运用口头和文字表达自己；善于充分调动集体和个人的积极性、创造性，协调各方面力量开展工作，高效完成任务；

（四） 原则上应当是本协会的成员；

（五） 最多可以连任两届；

（六） 经会员大会通过后须经登记管理机关审查后产生，若审查后被予以否定，不得当选。

\textbf{第三十一条：权限}

（一） 定期召开例会，制定协会的规章制度以及活动计划；

（二） 组织协会成员的联系和开展健康向上的活动；

（三） 做好工作记录、成员大会记录以及活动计划、活动总结；

（四） 培养、推荐协会的接班人，做好负责人的换届选举工作；

（五） 加强校内协会之间的交流，扩大本协会的影响，树立良好的协会形象，维护本协会的名誉，监督协会成员的学习与实践情况。

\textbf{第三十二条：罢免与辞职程序}

（一）罢免程序：

1. 协会负责人若严重失职或滥用职权，经活跃会员大会投票通过，可撤销其职务，报院团委审查；

2. 若协会负责人在任期内， 出现不符合担任协会干部的情形，院团委可直接予以撤职。

（二）辞职程序：

1. 协会负责人若需提前离职，应提前一个月提出，经秘书处备案，指导教师同意后，报院团委备案；

2. 递交辞呈给登记机关时，须提前一周通知，待新一届负责人选出后，方可退出职务，在职期间继续履行负责人的权利和义务；

3. 辞职者有责任按照程序做好交接工作。

\section{协会终止}

\textbf{第三十三条：终止动议提出}

本协会完成宗旨或自行解散或由于分立、合并等原因需要注销的， 由主席团提出终止动议。

\textbf{第三十四条：终止动议审批}

终止动议须经活跃会员大会表决通过，并报学校有关单位审查同意。

\textbf{第三十五条：清算程序}

协会终止前，须在人工智能学院及有关机关指导下成立清算组织，清理债权债务，处理善后事宜。清算期间，不开展清算以外的活动。

\textbf{第三十六条：注销登记}

协会经社团管理中心批准注销后即为终止。

\section{章程的修改}

\textbf{第三十七条：修改发起与筹备}

章程的修改起草工作由主席团负责。章程如需修改，需由主席团或五人以上发起，提交有关修改意见，交至主席团。

\textbf{第三十八条：修改表决方式}

修改章程应当 由活跃会员大会表决通过。活跃会员大会表决修改章程的决定，以全体活跃会员的三分之二以上出席，以出席人员三分之二以上通过。

\textbf{第三十九条：生效与公示}

修改决定确认通过后， 由秘书处对章程进行修改， 向全体会员（包括部员与非部员）公示三天后正式生效。

\textbf{第四十条：最终解释权}

本章程之最终解释权，由主席团先行组织集体商议、形成初步意见后，再依规报请学院审议核定。

2025年12月17日

\end{document}
